%==============================================================================
%  CAPÍTULO XIX
%==============================================================================
\chapter{Manual de tramitación: liquidación simplificada}
\label{cap:tram-liq-simplificada}

\fichaProcedimiento
  {Persona Deudora y micro o pequeña empresa.}
  {Juzgado civil del domicilio del deudor.}
  {De seis meses a un año, según existan o no bienes realizables.}
  {Reducido: sin martillero tradicional y con supresión de trámites.}
  {Extinción de los saldos insolutos, salvo incidente de mala fe o deudas excluidas.}

El desarrollo dogmático se encuentra en el capítulo
\ref{cap:liquidacion-simplificada}.

%==============================================================================
\bloqueEncargo
%==============================================================================

Este es el procedimiento de mayor volumen del sistema y aquel en que la asesoría produce
la diferencia más nítida entre un resultado y su opuesto. La razón es simple: el
beneficio que persigue el cliente ---la extinción de sus saldos insolutos--- puede
denegarse íntegramente por una omisión en el inventario.

\begin{foco}[Todo el encargo se juega en el inventario]
El procedimiento es sencillo. Lo que no es sencillo es declarar completamente. La
diferencia entre un \discharge{} obtenido y uno denegado suele reducirse a si el
inventario incluyó el APV, el vehículo que ya no se usa o el inmueble transferido al
cónyuge. El trabajo del abogado consiste, esencialmente, en asegurar que nada quede
fuera.
\end{foco}

%==============================================================================
\bloqueEntrevista
%==============================================================================

Sobre el cuestionario general del capítulo \ref{cap:entrevista}, este procedimiento exige
una indagación patrimonial exhaustiva. Las preguntas siguientes deben formularse de forma
específica, porque la pregunta genérica «¿tiene bienes?» produce respuestas incompletas
de buena fe.

\begin{cuestionario}[Barrido patrimonial específico]
  \pregunta{¿Tiene APV, cuenta 2, depósitos a plazo, fondos mutuos o acciones?}
  {Es la omisión más citada por la jurisprudencia como indicio de mala fe. El cliente no
  los considera «bienes».}

  \pregunta{¿Hay algún vehículo inscrito a su nombre, aunque no lo use o lo haya vendido
  sin transferir?}
  {La inscripción registral es lo que importa, no la posesión material. La «venta con
  poder» sigue apareciendo a su nombre.}

  \pregunta{¿Tiene alguna cuenta bancaria, aunque esté sin movimiento?}
  {Incluidas cuentas RUT y cuentas de ahorro olvidadas.}

  \pregunta{¿Es heredero de alguien, o tiene una herencia pendiente de partición?}
  {Los derechos hereditarios integran el patrimonio aunque no se hayan materializado.}

  \pregunta{¿Le deben dinero? ¿Tiene juicios en que sea demandante?}
  {Los créditos contra terceros son activo realizable y deben declararse.}

  \pregunta{¿Tiene participación en alguna sociedad, aunque esté inactiva?}
  {Los derechos sociales son bienes, aun cuando la sociedad no opere.}

  \pregunta{¿Ha recibido devoluciones de impuestos o las espera?}
  {Ingresan a la masa si se devengan durante el procedimiento.}
\end{cuestionario}

\begin{cuestionario}[Remuneración]
  \pregunta{¿Cuál es su remuneración bruta y líquida?}
  {Define si existe tramo excedente sujeto a incautación (capítulo
  \ref{cap:liquidacion-simplificada}).}

  \pregunta{¿Qué descuentos voluntarios tiene: APV, seguros, cuotas sociales?}
  {Los descuentos voluntarios reciben tratamiento distinto de los obligatorios en el
  cálculo del excedente.}

  \pregunta{¿Recibe bonos, gratificaciones o comisiones variables?}
  {Pueden situar meses puntuales por sobre el tramo inembargable.}
\end{cuestionario}

\begin{cuestionario}[Vivienda y entorno familiar]
  \pregunta{¿Dónde vive y a nombre de quién está la vivienda?}
  {La entrega voluntaria reemplaza a la incautación física, pero el cliente necesita saber
  con precisión qué ocurrirá en su casa.}

  \pregunta{¿Hay bienes en su hogar que pertenezcan a otras personas?}
  {Deben identificarse y documentarse antes, para excluirlos de la entrega.}

  \pregunta{¿Su cónyuge o conviviente tiene bienes propios? ¿Cuál es el régimen
  patrimonial del matrimonio?}
  {La sociedad conyugal altera sustancialmente la composición de la masa.}
\end{cuestionario}

%==============================================================================
\bloqueDocumental
%==============================================================================

\begin{checklist}[Barrido documental completo]
  \item Declaración jurada de antecedentes fidedignos conforme a la \ncg{22}
  \parencite{ncg-22}.
  \item Certificado de saldos de APV y cuenta 2 de la AFP y de aseguradoras.
  \item Cartolas de todas las cuentas bancarias de los últimos doce meses.
  \item Certificado de anotaciones vigentes del Registro Civil (barrido por RUT, no por
  vehículo conocido).
  \item Certificados de dominio vigente, hipotecas y gravámenes de cada inmueble.
  \item Certificado de matrimonio con anotaciones marginales, para acreditar el régimen
  patrimonial.
  \item Liquidaciones de remuneraciones de los últimos seis meses.
  \item Consulta de causas por RUT en el portal del Poder Judicial, como demandante y
  como demandado.
  \item Carpeta tributaria, para detectar devoluciones y participaciones societarias.
\end{checklist}

\begin{alerta}[Barrido por RUT, no por memoria]
Los certificados deben pedirse por RUT y no bien por bien. El objetivo es descubrir lo
que el cliente no recuerda tener, y eso sólo se consigue consultando los registros de
forma general.
\end{alerta}

%==============================================================================
\bloqueFocos
%==============================================================================

\subsection{La integridad del inventario}

Es el foco único y dominante. Toda la energía del encargo debe dirigirse a que el
inventario sea completo, porque de ello depende el resultado. Un activo de escaso valor
declarado no perjudica; el mismo activo omitido puede costar el \discharge{} entero.

\subsection{La entrega voluntaria}

La sustitución de la incautación física por la entrega voluntaria es la principal
humanización del procedimiento, pero está condicionada: la orden expresa del tribunal, la
infracción de los deberes de depositario o la detección de activos no declarados
reactivan la incautación. Es otro incentivo para declararlo todo.

\subsection{El tratamiento de la remuneración}

El tramo inembargable y el cálculo del excedente son la materia que más consultas genera
del cliente y más controversias con el empleador. Conviene anticiparse: informar al
cliente el cálculo estimado y, si hay retención, coordinar tempranamente con el
liquidador.

\subsection{Las deudas que sobrevivirán}

CAE, crédito solidario y pensiones de alimentos. El cliente debe tener claridad sobre
ellas desde el primer día y por escrito.

%==============================================================================
\bloqueCronograma
%==============================================================================

\begin{checklist}[Secuencia de trabajo]
  \item \textbf{Semanas $-4$ a $-1$.} Barrido documental completo por RUT. Construcción
  del inventario. Acta de diagnóstico firmada.
  \item \textbf{Día 0.} Presentación con declaración jurada de antecedentes fidedignos.
  \item \textbf{Resolución de liquidación.} Publicación en el \boletin{}. Nominación del
  liquidador.
  \item \textbf{Entrega de bienes.} Coordinación de la entrega voluntaria y exclusión
  documentada de bienes de terceros.
  \item \textbf{Verificación de créditos.} Control de las verificaciones y objeción de
  las improcedentes.
  \item \textbf{Realización.} Subasta electrónica de los bienes realizables.
  \item \textbf{Resolución de término.} Extinción de los saldos insolutos, salvo
  incidente de mala fe o deudas excluidas.
\end{checklist}

%==============================================================================
\bloqueErrores
%==============================================================================

\errorfrecuente{Preguntar «¿tiene bienes?» en lugar de hacer el barrido específico}
{El cliente responde de buena fe que no, omitiendo APV, vehículo inscrito o derechos
sociales, y se configura la causal del \art{169 A} \Nro 1.}
{Recorrer el cuestionario de barrido patrimonial pregunta por pregunta y contrastarlo con
certificados obtenidos por RUT.}

\errorfrecuente{Confiar en que el cliente sabe qué vehículos tiene inscritos}
{Aparece un vehículo «vendido con poder» que sigue a su nombre, y se lee como
ocultamiento.}
{Solicitar el certificado de anotaciones vigentes por RUT.}

\errorfrecuente{No verificar el régimen patrimonial del matrimonio}
{Se declara mal la composición de la masa y surgen conflictos con el cónyuge durante la
entrega.}
{Pedir siempre el certificado de matrimonio con anotaciones marginales.}

\errorfrecuente{No preparar al cliente para la entrega de bienes}
{La diligencia se convierte en un conflicto doméstico y, si hay resistencia, puede
reactivarse la incautación forzada.}
{Explicar el procedimiento con antelación e identificar por escrito los bienes de
terceros.}

\errorfrecuente{Prometer la extinción de todas las deudas}
{El cliente termina con el CAE o los alimentos intactos y con la convicción de haber sido
engañado.}
{Consignar por escrito, en el acta de diagnóstico y en el mandato, qué deudas
probablemente sobrevivirán.}
